\documentclass[UTF8,letter,10pt]{ctexart}
\usepackage{geometry}
\geometry{left=2.5cm, right=2.5cm, top=2.5cm, bottom=2.5cm}
\setmainfont{Times New Roman}
%\setmainfont{TeX Gyre Termes}
\usepackage{booktabs}
\usepackage{colortbl}
\usepackage{xcolor}
\usepackage{microtype}
\usepackage{multirow}
\usepackage{makecell}
\usepackage{graphicx}
\usepackage{float}
\usepackage{hyperref}
\usepackage{amssymb}

\definecolor{intjpurple}{HTML}{88619A}
%\definecolor{intjpurple}{HTML}{4A148C}

\newcommand{\verified}{\colorbox{green}{\textcolor{white}{\checkmark Verified}}}


\begin{document}

\listoftables

\begin{table}[h!]
    \caption[八功能关键词 \verified]{八功能关键词 | Semantic Field of Each Function-Attitude.
    The middle word is the apex of a triangle.
    The first of the three words (as we read across for each of the eight mental processes)
    can be thought of as what the process looks like on the surface, to another person seeing someone for the first time using that process -- what we might call its \textit{persona} level.
    The second, central, word captures the heart of the process -- the one that is
    embraced by the \textit{ego} as its chief concern.
    The third word embodies what the process becomes at its most evolved level when it
    is working in sympathetic harmony with what Jung call the \textit{Self}:
    one might call this the goal of the process.\\
    \textcolor{gray}{\small Source: Page 4, Section ``The Eight Function-Attitudes Unpacked'' from Beebe~\cite{beebe}.}}
    \begin{center}

        \begin{tabular}{ccccc}
            \toprule
            \multirow{2}{*}{\textbf{Function-attitude}} & \multirow{2}{*}{\textbf{Full Name}} & \multicolumn{3}{c}{\textbf{Semantic Field Triangle}} \tabularnewline
            \cmidrule(lr){3-5}
            & & \small{面具 (Persona)} & \small{自我 (Ego)} & \small{自性 (Self)} \tabularnewline
            \midrule

            \rowcolor{black!10}Se & Extraverted Sensing & Engaging & Experiencing & Enjoying \tabularnewline
            Si & Introverted Sensing & Implementing & Verifying & Accounting \tabularnewline

            \rowcolor{black!10}Ne & Extraverted Intuition & Entertaining & Envisioning & Enabling \tabularnewline
            Ni & Introverted Intuition & Imagining & Knowing & Divining \tabularnewline

            \rowcolor{black!10}Te & Extraverted Thinking & Regulating & Planning & Enforcing \tabularnewline
            Ti & Introverted Thinking & Naming & Defining & Understanding \tabularnewline

            \rowcolor{black!10}Fe & Extraverted Feeling & Validating & Affirming & Relating \tabularnewline
            Fi & Introverted Feeling & Judging & Appraising & Establishing the Value \tabularnewline

            \bottomrule
        \end{tabular}

    \end{center}
\end{table}


\begin{table}[h!]
\caption[MBTI - 荣格八维映射表 \verified]{MBTI - 荣格八维映射表 (Function Stacks)}
    \label{tab:mbti-jung}
    \centering
    \begin{tabular}{cccc}
        \toprule
        \multicolumn{4}{c}{\textbf{NT -- 分析师 Analysts}} \\
        \cellcolor{intjpurple}\color{white} INTJ (Ni-Te-Fi-Se) & INTP (Ti-Ne-Si-Fe) & ENTJ (Te-Ni-Se-Fi) & ENTP (Ne-Ti-Fe-Si) \\
        \midrule
        \multicolumn{4}{c}{\textbf{NF -- 外交官 Diplomats}} \\
        INFJ (Ni-Fe-Ti-Se) & INFP (Fi-Ne-Si-Te) & ENFJ (Fe-Ni-Se-Ti) & ENFP (Ne-Fi-Te-Si) \\
        \midrule
        \multicolumn{4}{c}{\textbf{SJ -- 守护者 Sentinels}} \\
        ISTJ (Si-Te-Fi-Ne) & ISFJ (Si-Fe-Ti-Ne) & ESTJ (Te-Si-Ne-Fi) & ESFJ (Fe-Si-Ne-Ti) \\
        \midrule
        \multicolumn{4}{c}{\textbf{SP -- 探索者 Explorers}} \\
        ISTP (Ti-Se-Ni-Fe) & ISFP (Fi-Se-Ni-Te) & ESTP (Se-Ti-Fe-Ni) & ESFP (Se-Fi-Te-Ni) \\
        \bottomrule
    \end{tabular}
\end{table}

\begin{table}[h!]
\caption{MBTI 类型总览表 (功能栈、原型、核心特质与代表名人)}
\resizebox{\textwidth}{!}{%
\begin{tabular}{@{}l l c l l@{}}
\toprule
\textbf{类型} & \textbf{功能栈 (Stack)} & \textbf{原型 (Archetype)} & \textbf{核心特质 (Key Traits)} & \textbf{代表名人 (Notable Figures)} \\
\midrule
\multicolumn{5}{c}{\textbf{分析家 (Analysts) - NT (理性与筹划)}} \\
\midrule
\cellcolor{intjpurple}\color{white}\textbf{INTJ} & Ni - Te - Fi - Se & 建筑师 & 战略宏大，系统构建，独立深邃 & 牛顿, 尼采, 马克思 \\
\textbf{INTP} & Ti - Ne - Si - Fe & 逻辑学家 & 追求真理，逻辑自洽，抽象解构 & 爱因斯坦, 达尔文, 林肯 \\
\textbf{ENTJ} & Te - Ni - Se - Fi & 指挥官 & 杀伐决断，整合资源，极强执行 & 拿破仑, 撒切尔, 凯撒 \\
\textbf{ENTP} & Ne - Ti - Fe - Si & 辩论家 & 思维跳跃，挑战传统，从零到一 & 达芬奇, 苏格拉底, 富兰克林 \\
\midrule
\multicolumn{5}{c}{\textbf{外交家 (Diplomats) - NF (理想与共情)}} \\
\midrule
\textbf{INFJ} & Ni - Fe - Ti - Se & 提倡者 & 洞察人性，精神探索，利他主义 & 荣格, 柏拉图, 甘地 \\
\textbf{INFP} & Fi - Ne - Si - Te & 调停者 & 忠于自我，内心丰富，悲悯治愈 & 莎士比亚, 托尔金, 梵高 \\
\textbf{ENFJ} & Fe - Ni - Se - Ti & 主人公 & 极具感召，洞察群体，愿景激励 & 马丁·路德·金, 歌德 \\
\textbf{ENFP} & Ne - Fi - Te - Si & 竞选者 & 热情洋溢，追求自由，激发灵感 & 马克·吐温, 迪士尼, 王尔德 \\
\midrule
\multicolumn{5}{c}{\textbf{守护者 (Sentinels) - SJ (秩序与责任)}} \\
\midrule
\textbf{ISTJ} & Si - Te - Fi - Ne & 物流师 & 恪尽职守，维护规则，实事求是 & 华盛顿, 奥古斯都, 贝索斯 \\
\textbf{ISFJ} & Si - Fe - Ti - Ne & 守卫者 & 细腻体贴，幕后奉献，守护和谐 & 特蕾莎修女, 罗莎·帕克斯 \\
\textbf{ESTJ} & Te - Si - Ne - Fi & 总经理 & 高效务实，建立制度，崇尚秩序 & 福特, 洛克菲勒, 杜鲁门 \\
\textbf{ESFJ} & Fe - Si - Ne - Ti & 执政官 & 热情好客，重视地位，服务群体 & 卡内基, 比尔·克林顿 \\
\midrule
\multicolumn{5}{c}{\textbf{探险家 (Explorers) - SP (行动与当下)}} \\
\midrule
\textbf{ISTP} & Ti - Se - Ni - Fe & 鉴赏家 & 冷静克制，精通技艺，危机处理 & 宫本武藏, 隆美尔, 李小龙 \\
\textbf{ISFP} & Fi - Se - Ni - Te & 探险家 & 审美独特，活在当下，打破常规 & 迈克尔·杰克逊, 莫扎特 \\
\textbf{ESTP} & Se - Ti - Fe - Ni & 企业家 & 敏锐机警，利益导向，冒险犯难 & 丘吉尔, 海明威, 亚历山大 \\
\textbf{ESFP} & Se - Fi - Te - Ni & 表演者 & 表现力强，乐享生活，感染力强 & 玛丽莲·梦露, 猫王 \\
\bottomrule
\end{tabular}}

算法：MBTI 到认知功能栈的转换（简易版）
\begin{enumerate}
        \footnotesize
    \item \textbf{第一步：J/P 指针（谁获得下标 $e$？）} \\
    观察类型的\textbf{最后一个字母}。它决定了哪个功能是\textbf{外倾的 ($X_e$)}。
    \begin{itemize}
        \item 如果是 \textbf{J}：\textbf{判断}功能（$T$ 或 $F$）变为外倾（$Te$ 或 $Fe$）。
        \item 如果是 \textbf{P}：\textbf{感知}功能（$S$ 或 $N$）变为外倾（$Se$ 或 $Ne$）。
        \item \textit{注：剩余的那个功能自动变为内倾（$X_i$）。}
        \item \textit{速记：\textbf{iP}(助记：IP 地址) 和 \textbf{eJ}，即和最后一个字母 \textbf{J/P} 直接相连的 \textbf{T/F} 的下标，和最后一个字母的组合一定是 \textbf{iP} 或 \textbf{eJ}。}
    \end{itemize}

    \item \textbf{第二步：E/I 指针（谁是主导？）} \\
    观察类型的\textbf{第一个字母}。它决定了\textbf{主导功能}的位置。
    \begin{itemize}
        \item 如果是 \textbf{E}：带有 \textbf{$e$} 下标的功能排在第一位。
        \item 如果是 \textbf{I}：带有 \textbf{$i$} 下标的功能排在第一位。
    \end{itemize}

    \item \textbf{第三步：镜像补全（推导第3与第4功能）} \\
    功能栈的后半部分仅仅是前半部分的\textbf{反向镜像}。
    \begin{itemize}
        \item \textbf{第3功能} = 第2功能的完全对立面（如 $Te \to Fi$）。
        \item \textbf{第4功能} = 第1功能的完全对立面（如 $Ni \to Se$）。
    \end{itemize}

    \item \textbf{第四步：脚标反转补全（推导第5,6,7,8功能）} \\
    影子栈(后4位)的是前4位功能栈的\textbf{第一个大写字母相同，第二个小写字母反转}。
    E.g. \textbf{第5功能} = 第1功能的大写字母相同，小写字母反转（如 $Ni \to Ne$）；
    \textbf{第6功能} = 第2功能的大写字母相同，小写字母反转（如 $Te \to Ti$）；
    \textbf{第7功能} = 第3功能的大写字母相同，小写字母反转（如 $Fi \to Fe$）；
    \textbf{第8功能} = 第4功能的大写字母相同，小写字母反转（如 $Se \to Si$）。
\end{enumerate}

\end{table}


\begin{table}[h!]
    \caption[认知功能的1*2*标准健康表现]{认知功能的标准健康表现(The Ideal State)。
这通常是1*主导(Hero)与2*辅助(Parent)功能的天然状态。
在个体化(Individuation)的高阶整合中，即便是劣势功能，通过刻意练习也能呈现出此表的特质。
\textbf{核心区别}：原生功能是自动化的“充电器”(Charging)，而通过后天训练整合的功能——尽管外部表现完美——本质上仍是需要消耗专注力的“耗电器”(Draining)。}
    \label{tab:cog-function-normal}
    \begin{center}
\begin{tabular}{cccc}
\toprule
简称 & 全称 & 核心词 & 正常表现\tabularnewline
\midrule
Se & Extraverted Sensing & 体验，当下，刺激 & 活在当下，直接行动，现实主义\tabularnewline
Si & Introverted Sensing & 经验，细节，稳健 & 参照经验，细节敏感，追求守序\tabularnewline
Ne & Extraverted Intuition & 发散，脑洞，可能 & 跳跃思维，探索新奇，抽象关联\tabularnewline
Ni & Introverted Intuition & 收敛，洞察，本质 & 预测未来，现象本质，单一愿景\tabularnewline
Te & Extraverted Thinking & 效率，执行，调配 & 结果导向，逻辑外化，强势直接\tabularnewline
Ti & Introverted Thinking & 逻辑，解构，精准 & 追求真理，刨根问底，无视权威\tabularnewline
Fe & Extraverted Feeling & 和谐，人际，共情 & 情绪传染，维护关系，社交礼仪\tabularnewline
Fi & Introverted Feeling & 本心，真实，价值 & 忠于自我，情感内化，排斥虚伪\tabularnewline
\bottomrule
\end{tabular}
    \end{center}


\begin{description}
        \footnotesize

    \item[Se (外倾感觉 | Extraverted Sensing) —— \textit{当下的冲击}] \hfill \\
    它是意识与物理现实之间\textbf{零延迟}的撞击。Se 不仅仅是看和听，它是\textbf{拒绝抽象}，是渴望被客观世界的\textbf{原始数据}直接穿透的冲动。它是一场感官的贪婪盛宴，追求的是\textbf{“在场”的极致强度}。对于Se而言，真理不在书本里，而在此时此刻皮肤的触感、发动机的轰鸣和色彩的饱和度中；它是一种\textbf{将自我完全溶解在客观情境}中的能力，不仅是观察者，更是这一秒钟现实的共舞者。

    \item[Si (内倾感觉 | Introverted Sensing) —— \textit{时间的琥珀}] \hfill \\
    它是客观世界在主观灵魂中留下的\textbf{永久压痕}。Si 并不记录客观事实，它记录的是\textbf{事实引发的主观回响}。它是一种将流动的现实凝固为\textbf{原型（Archetype）}的能力，将当下的瞬间不断与内心的“档案库”进行比对。它追求的是\textbf{恒常性}，像守护神一样站在时间的长河中，通过仪式、习惯和身体记忆，抵御外界的混乱与无常，构建出一个充满历史厚度与个人神话的安全堡垒。

    \item[Ne (外倾直觉 | Extraverted Intuition) —— \textit{可能性的超新星}] \hfill \\
    它是对\textbf{“现状”的永恒不满}，以及对\textbf{“潜在”的无限饥渴}。Ne 透过现象看到的是背后无形的连接线，它是一种\textbf{发散性的创造力}，永远在寻找下一个跳板。对于Ne来说，一个物体不是它本身，而是它\textbf{“可能变成的一切”}。它是一种\textbf{熵增}的认知模式，拒绝结构的封闭，热衷于在看似无关的事物间引爆火花，通过不断的重组与联想，试图穷尽宇宙中所有的可能性组合。

    \item[Ni (内倾直觉 | Introverted Intuition) —— \textit{深渊的凝视}] \hfill \\
    它是意识从纷繁的表象中\textbf{撤退}，转向潜意识深处寻找\textbf{单一核心模式}的过程。Ni 是大脑的\textbf{后台压缩算法}，它忽略细节，只提取那个跨越时间的\textbf{元叙事}。它是一种\textbf{收敛性的洞察}，往往以意象、幻觉或预感的面目出现，不讲逻辑，直接呈现结论。它追求的是\textbf{终极的确定性}，试图透过时间的迷雾，用一只眼睛看到事物发展的终点，是一种带有宿命论色彩的、指向未来的透视眼。

    \item[Te (外倾思维 | Extraverted Thinking) —— \textit{秩序的建筑师}] \hfill \\
    它是将内在逻辑\textbf{强加于}混乱世界之上的意志力。Te 不在乎原本的真理是什么，只在乎\textbf{什么行得通}（Workable）。它是一种\textbf{以结果为导向的结构化思维}，将万物（包括人）视为资源、参数和变量，通过制定规则、标准和流程来征服熵。它是\textbf{客观逻辑的利剑}，也是社会运转的骨架，其本质是对\textbf{效率}和\textbf{可控性}的极致追求，任何无法被度量的事物在Te眼中都缺乏存在的合法性。

    \item[Ti (内倾思维 | Introverted Thinking) —— \textit{真理的解剖刀}] \hfill \\
    它是对\textbf{概念纯度}和\textbf{逻辑自洽}的绝对洁癖。Ti 并不关心外界是否认可，只关心思维大厦内部的\textbf{结构是否完美无瑕}。它是一种\textbf{主观的本质主义}，试图剥离一切冗余的语境，给事物重新命名，直到找到那个不可再分的定义。对于Ti而言，现实世界往往是充满了错误的，只有那个在脑海中构建的、如水晶般剔透的逻辑模型，才是\textbf{永恒的真理}。

    \item[Fe (外倾情感 | Extraverted Feeling) —— \textit{共鸣的指挥家}] \hfill \\
    它是将\textbf{自我边界消融}以融入集体情感流动的本能。Fe 不是个人的情绪，而是\textbf{客观的情感价值}。它是一种\textbf{动态的社交调节器}，时刻扫描环境中的人际氛围，并本能地调整自己的频率以达成\textbf{和谐}。它追求的是\textbf{人与人之间的连接感}，通过确认共同的价值观、礼仪和道德规范，将分散的个体凝聚为一个整体。在Fe眼中，\textbf{“我们”}的存在先于且高于\textbf{“我”}的存在。

    \item[Fi (内倾情感 | Introverted Feeling) —— \textit{静默的圣火}] \hfill \\
    它是对\textbf{个体情感深度}和\textbf{内在价值}的誓死捍卫。Fi 是一口深井，水面波澜不惊，井底却涌动着\textbf{炽热的岩浆}。它拒绝被外界的道德标准同化，只听从\textbf{内心的良知}与\textbf{真实的喜恶}。它是一种\textbf{排他性的评价体系}，追求的是\textbf{绝对的真诚（Authenticity）}，为了守护内心那片神圣不可侵犯的净土，它宁愿与全世界为敌，是一种虽无声却极具穿透力的精神力量。

\end{description}

\end{table}


\begin{table}[h!]
    \caption[认知功能的角色表]{认知功能的角色表（含标准定义符合度）}
    \label{tab:cog-function-roles}
    \resizebox{\textwidth}{!}{%
    \begin{tabular}{ccccc}
    \toprule
    序号 & 位置 & 角色名称 & 典型表现 & 对标“正常表现”程度\tabularnewline
    \midrule
        1 & 主导 (Dominating) & 英雄(Hero) & 核心力量 & 120\% (本能超载：如呼吸般自然)\tabularnewline
        2 & 辅助 (Auxiliary) & 家长(Parent) & 辅助工具 & 100\% (完美对标：最成熟可控)\tabularnewline
        3 & 第三 (Tertiary) & 孩子(Child) & 舒适区 & 50\% $\sim$ 70\% (任性/游戏：表现不稳定)\tabularnewline
        4 & 劣势 (Inferior) & 阿尼玛(Anima) & 压力源泉/理想 & 10\% (原始) $\rightarrow$ 70\% (整合后工具化)\tabularnewline
    \midrule
        5 & 影子 & 反向作用(Opposing) & 质疑者 & -20\% (抗拒/相反：态度厌恶)\tabularnewline
        6 & 影子 & 批判家长(Critical) & 挑刺者 & -50\% (双标/虚伪：严于律人)\tabularnewline
    \rowcolor{gray!17}7 & 影子 & 盲点(Trickster) & 真正盲区 & 0\% (黑洞/混乱：无法理解)\tabularnewline
        8 & 影子 & 恶魔(Demon) & 毁灭者 & -100\% (彻底颠覆：旨在摧毁)\tabularnewline
    \bottomrule
    \end{tabular}}

    \begin{description}
            \footnotesize

    \item[1. 主导功能 (Dominant / The Hero) —— \textit{意识的绝对王座}] \hfill \\
    它是自我（Ego）的\textbf{同义词}，是你感知世界的出厂设定透镜。它过于强大、过于熟练，以至于你感觉不到它的存在——因为\textbf{它就是“你”本身}。它是人生叙事中的\textbf{英雄主角}，拥有不受质疑的统治权。它不仅定义了你的擅长，更定义了你的\textbf{傲慢}。它像正午的太阳，虽然照亮了世界，但强烈的光芒也制造了最深沉的阴影；它让人产生一种全能的幻觉，认为世界本该如此运行，却不知自己已被锁死在单一维度的现实中。

    \item[2. 辅助功能 (Auxiliary / The Parent) —— \textit{责任的守望者}] \hfill \\
    它是英雄身边忠诚的\textbf{辅佐者}，也是连接内心与外界的\textbf{桥梁}。不同于主导功能的本能与冲动，辅助功能总是带有一种\textbf{“养育”与“责任”}的色彩。它是你用来帮助他人、回馈社会的工具，因此显得成熟、稳重且富有建设性。它是防止英雄陷入孤芳自赏的\textbf{刹车片}，通过提供相反维度的视角（内倾者提供外倾视角，反之亦然），迫使自我面对现实的复杂性，将英雄的能量转化为可被社会接纳的成果。

    \item[3. 第三功能 (Tertiary / The Child) —— \textit{诱惑的避风港}] \hfill \\
    它是永恒少年（Puer Aeternus）的游乐场，是自我最喜欢的\textbf{舒适区}。它不负责任、充满玩心、天真烂漫，同时也脆弱易碎。它是你\textbf{放松防御}的地方，也是你用来自我安抚的奶嘴。然而，它极具\textbf{欺骗性}，因为它常常绕过严厉的“父母功能”，与主导功能结成\textbf{“恶性循环（Loop）”}的同盟，诱导你退回幼稚的自我中心状态，用一种看似无害的方式合理化你的逃避与停滞。

    \item[4. 劣势功能 (Inferior / The Anima) —— \textit{黄昏的守门人}] \hfill \\
    它是意识与潜意识的接驳点，是\textbf{阿喀琉斯之踵}。它是所有抱负的起点，也是所有崩溃的终点。作为那个总是被我们笨拙模仿的“理想自我”，它承载了极度的羞耻感与极度的神圣感。它是通往完整的必经之路，要求英雄必须经历\textbf{“自我羞辱”}般的痛苦，承认自己的无能，才能从那个一直被忽视的角落里，捡起遗失灵魂的另一半。

    \item[5. 第五功能 (Opposing / The Nemesis) —— \textit{怀疑的暗影}] \hfill \\
    它是主导功能的\textbf{负面镜像}，是潜意识阴影中的\textbf{反对党领袖}。它时刻都在通过\textbf{偏执与怀疑}来挑战英雄的权威。每当你感到莫名的被针对、或者对他人的动机产生毫无根据的敌意时，就是它在低语。它的存在是为了提醒英雄：你的视角不仅是有限的，甚至可能是错误的。它是一面\textbf{抗拒之墙}，通过顽固和被动攻击，试图阻碍自我的单边扩张，它是你最难缠的假想敌。

    \item[6. 第六功能 (Critical Parent / The Witch or Senex) —— \textit{羞耻的审判官}] \hfill \\
    它是辅助功能的阴暗面，是一个拥有毁灭性打击力的\textbf{内部批评家}。它不像“父母功能”那样去养育，而是利用它对规则的深刻洞察去\textbf{贬低、限制和羞辱}。当它指向他人时，你会变得刻薄、说教且令人窒息；当它指向自己时，它是那个永远让你觉得自己“不够好、不够格、是个失败者”的声音。它掌管着\textbf{“瘫痪”}的力量，能精准地找到软肋，用一句话就让自我的努力化为乌有。

    \item[7. 第七功能 (Trickster / The Blindspot) —— \textit{荒谬的哈哈镜}] \hfill \\
    它是心理结构中的\textbf{黑洞}，是你完全\textbf{看不见}也\textbf{搞不懂}的盲区。它不是简单的弱点，而是\textbf{概念的缺失}。在这里，你是滑稽的、被双重束缚的、被命运捉弄的小丑。它负责制造生活中的\textbf{荒诞剧}，通过让你反复跌入同一个陷阱、犯下低级且令人啼笑皆非的错误，来嘲弄自我的严肃性。它的残酷在于它剥夺了你的认知工具，让你在这一领域赤身裸体，不得不面对\textbf{混沌}本身。

    \item[8. 第八功能 (Demon / The Daemon) —— \textit{虚无的毁灭者}] \hfill \\
    它是潜意识最深处的\textbf{核按钮}，是人格结构的基岩。它是劣势功能的阴暗面，代表着\textbf{绝对的陌生}与\textbf{非人（Inhuman）}的特质。平时它处于休眠状态，但当自尊被彻底践踏、自我无路可走时，它会苏醒。它不寻求解决问题，只寻求\textbf{彻底的破坏}——毁灭那个造成痛苦的世界，或者毁灭你自己。然而，在神秘主义视角下，这种彻底的虚无也是\textbf{涅槃}的前夜，只有当旧的人格结构被它完全粉碎，新的神性（Daimonic）才可能从中升起。

\end{description}
\end{table}



\begin{table}[h!]
    \caption[认知功能的3*永恒之子(Child)表现]{认知功能的3*永恒之子(Child)表现： 治愈、游戏与任性。这是人格的“游乐场”与退行区。表现形式：天真单纯的创造力(Playfulness)，或是压力下像被宠坏的孩子般无理取闹(Immaturity)。}
    \label{tab:cog-function-child}
    \centering
    \resizebox{\textwidth}{!}{
\begin{tabular}{ccccc}
\toprule
孩童功能 & 涉及类型 & 核心台词 & 游戏表现 (Playfulness) & 幼稚表现 (Immaturity)\tabularnewline
\midrule
\multicolumn{5}{c}{理性游戏组：Judgement}\tabularnewline
\midrule
Ti  & INFJ, ISFJ & 居然是这样！ & 纯粹好奇，钻研冷门，解谜乐趣 & 逻辑死板，钻牛角尖，固执己见\tabularnewline
Fi  & \cellcolor{intjpurple}\color{white}INTJ, ISTJ & 我偏要喜欢 & 傲娇(Tsundere)，深情忠诚，情感纯粹 & 情感任性，无视大局，敏感记仇\tabularnewline
Te  & ENFP, ESFP & 我能搞定！ & 虚张声势，突发指挥，热衷整理 & 盲目自信，拒绝援助，三分钟热度\tabularnewline
Fe  & ENTP, ESTP & 大家都爱我 & 迷人捣蛋，幽默风趣，活跃气氛 & 渴望关注，哗众取宠，惧怕被弃\tabularnewline
\midrule
\multicolumn{5}{c}{感知游戏组：Perception}\tabularnewline
\midrule
Ni  & ISTP, ISFP & 我有预感... & 迷信直觉，神秘主义，宿命感 & 盲信猜测，无端幻想，消极宿命\tabularnewline
Si  & INTP, INFP & 还是旧的好 & 恋旧收藏，沉溺舒适，重复安稳 & 抗拒改变，退行童年，逃避现实\tabularnewline
Ne  & ESTJ, ESFJ & 如果这样呢？ & 突发脑洞，冷笑话，无厘头创意 & 容易分心，盲目乐观，无法落地\tabularnewline
Se  & ENTJ, ENFJ & 我全都要！ & 报复享乐，追求奢华，感官盛宴 & 挥霍无度，暴饮暴食，缺乏节制\tabularnewline
\bottomrule
\end{tabular}}
\end{table}


\begin{table}[h!]
    \caption[认知功能的4*阿尼玛(Inferior)表现：正常态]{认知功能的4*阿尼玛(Inferior)表现：正常态 (Aspiration)。这是人格的“天真孩童”模式。表现形式：在安全、低压环境下，我们会拙劣但快乐地使用劣势功能，表现出一种单纯的向往和治愈感。}
    \label{tab:cog-function-aspiration}
\centering
\resizebox{\textwidth}{!}{
\begin{tabular}{ccccc}
\toprule
劣势功能 & 涉及类型 & 潜意识心态 (Mental Stance) & 表现特征 (Playful/Healing) & 典型行为案例\tabularnewline
\midrule
\multicolumn{5}{c}{感知型阿尼玛：Perception}\tabularnewline
\midrule
Se & \cellcolor{intjpurple}\color{white}INTJ, INFJ & “我想体验真实的美” & 精致享乐，审美强迫 & 偶尔大餐，沉迷Hi-Fi/摄影，撸猫\tabularnewline
Si & ENTP, ENFP & “我想拥有根基” & 仪式感，片刻怀旧 & 收藏童年玩具，过分执着某个节日\tabularnewline
Ne & ISTJ, ISFJ & “我想变得有趣” & 冷幽默，安全冒险 & 讲双关语笑话，看科幻片，交怪朋友\tabularnewline
Ni & ESTP, ESFP & “我想看透本质” & 迷信玄学，瞬间深沉 & 相信星座风水，酒后探讨人生意义\tabularnewline
\midrule
\multicolumn{5}{c}{判断型阿尼玛：Judgement}\tabularnewline
\midrule
Te & INFP, ISFP & “我想变得干练” & 整理癖，清单控 & 买很多收纳工具，列做不到的计划表\tabularnewline
Ti & ENFJ, ESFJ & “我想拥有智慧” & 引经据典，崇拜专家 & 读艰深晦涩的书，引用名人名言\tabularnewline
Fe & INTP, ISTP & “我想被人喜欢” & 憨厚微笑，笨拙示好 & 极其客气，送直男礼物，默默帮忙\tabularnewline
Fi & ENTJ, ESTJ & “我也有温柔一面” & 隐秘深情，道德洁癖 & 溺爱宠物，对某首老歌流泪，讲义气\tabularnewline
\bottomrule
\end{tabular}}
\end{table}


\begin{table}[h!]
    \caption[认知功能的4*阿尼玛(Inferior)表现：压抑/禁欲态]{认知功能的4*阿尼玛(Inferior)表现：压抑/禁欲态 (Suppression/Asceticism)。这是人格的“防御性屏蔽”模式。表现形式：为了防止劣势功能的“噪音”干扰主导功能，通过极端的规则化、极简化或回避，将其“冷冻”或“锁死”。}
    \label{tab:cog-function-suppression}
\centering
\resizebox{\textwidth}{!}{
\begin{tabular}{ccccc}
\toprule
劣势功能 & 涉及类型 & 屏蔽逻辑 (Suppression Logic) & 表现特征 (Defensive Minimalism) & 典型行为案例\tabularnewline
\midrule
\multicolumn{5}{c}{感知型阿尼玛：Perception}\tabularnewline
\midrule
Se & \cellcolor{intjpurple}\color{white}INTJ, INFJ & “感官是Ni的噪音” & 苦行僧，感官剥离 & 全黑穿搭，极简装修，甚至忘记吃饭\tabularnewline
Si & ENTP, ENFP & “身体是Ne的累赘” & 肉体忽视，反常规 & 昼夜颠倒，生活区像垃圾场，拒绝体检\tabularnewline
Ne & ISTJ, ISFJ & “变数是Si的敌人” & 极端保守，拒绝想象 & 刻板重复昨日，不仅不做梦，还打击梦想\tabularnewline
Ni & ESTP, ESFP & “未来是Se的恐惧” & 激进务实，嘲讽深刻 & “想那么多没用”，只活在这一秒，拒绝规划\tabularnewline
\midrule
\multicolumn{5}{c}{判断型阿尼玛：Judgement}\tabularnewline
\midrule
Te & INFP, ISFP & “结构是Fi的牢笼” & 被动抵抗，混乱防御 & 故意拖延，拒绝任何时间表，以无序为荣\tabularnewline
Ti & ENFJ, ESFJ & “质疑是Fe的破坏” & 智力依赖，盲从共识 & “大家都这么说”，拒绝独立思考，屏蔽异见\tabularnewline
Fe & INTP, ISTP & “社交是Ti的消耗” & 社交隐身，情感断连 & 无论谁发信都不回，完全切断人际信号\tabularnewline
Fi & ENTJ, ESTJ & “情绪是Te的低效” & 情感切除，机器人化 & “只谈公事”，无视他人痛苦，否认自己有情绪\tabularnewline
\bottomrule
\end{tabular}}
\end{table}


\begin{table}[h!]
    \caption[认知功能的4*阿尼玛(Inferior)表现：劫持/抓取态]{认知功能的4*阿尼玛(Inferior)表现：劫持/抓取态 (The Grip)。这是人格的“至暗时刻”。表现形式：主导功能崩溃，潜意识中的“阴影”强行接管意识，不仅智商下线，还会变成自己曾经最讨厌的那种人的“低配版”。}
    \label{tab:cog-function-grip}
\centering
\resizebox{\textwidth}{!}{
\begin{tabular}{ccccc}
\toprule
劣势功能 & 涉及类型 & 崩溃扳机 (The Trigger) & 内在心态 (Internal State) & 外显行为 (External Behavior)\tabularnewline
\midrule
\multicolumn{5}{c}{感知型阿尼玛：Perception}\tabularnewline
\midrule
Se & \cellcolor{intjpurple}\color{white}INTJ, INFJ & 脑力耗竭，现实细节失控 & “这世界太烂了，我也烂透了” & 暴饮暴食，疯狂购物，沉迷色情/游戏，强迫性整理\tabularnewline
Si & ENTP, ENFP & 无限可能变成混乱，生病 & “一切都完了，永远好不了了” & 隧道视野，甚至感觉身体正在腐烂(疑病)，死抠细节\tabularnewline
Ne & ISTJ, ISFJ & 熟悉流程被打断，未知风险 & “灾难马上就要降临” & 灾难化想象(Catastrophizing)，看到到处都是阴谋\tabularnewline
Ni & ESTP, ESFP & 此路不通，未来迷茫 & “人生毫无意义，我是受害者” & 偏执妄想，把中性眼神解读为恶意，相信宿命论\tabularnewline
\midrule
\multicolumn{5}{c}{判断型阿尼玛：Judgement}\tabularnewline
\midrule
Te & INFP, ISFP & 价值观被践踏，感到无能 & “既然没人懂我，那我就毁了一切” & 暴君模式，极度苛刻批判他人，逻辑混乱的独裁\tabularnewline
Ti & ENFJ, ESFJ & 人际背叛，缺乏反馈 & “人类是愚蠢且无可救药的” & 冷酷的犬儒主义，用刺耳的逻辑攻击他人软肋\tabularnewline
Fe & INTP, ISTP & 逻辑失效，被群体排斥 & “没人爱我，我是多余的” & 情绪崩溃大哭，极其敏感，像个受伤的孩子\tabularnewline
Fi & ENTJ, ESTJ & 权威丧失，内鬼背叛 & “全世界都对不起我” & 陷入自怜自艾(Self-pity)，情绪化爆发，甚至抑郁\tabularnewline
\bottomrule
\end{tabular}}
\end{table}


\begin{table}[h!]
\caption[极端压力/抓取状态(Grip)的解药]{极端压力/抓取状态(Grip)的解药。当主导功能枯竭时，切忌（用主导功能）硬抗，此时唯一的解药是安抚劣势功能，顺应其最原始的需求。}
    \label{tab:cog-function-inferior-remedy}
    \centering
    \resizebox{\textwidth}{!}{
\begin{tabular}{ccccc}
\toprule
劣势功能 & 涉及类型 & 核心策略 (Strategy) & 停止行为 (Stop) & 执行清单 (Action)\tabularnewline
\midrule
\multicolumn{5}{c}{感知解药：回归生理与当下 (Perception)}\tabularnewline
\midrule
Se & \cellcolor{intjpurple}\color{white}INTJ, INFJ & \textbf{低压感官疗法} & 停止挖掘意义，停止死磕理论 & 高质量睡眠，接触大自然，摄影，无目的散步\tabularnewline
Si & ENTP, ENFP & \textbf{熟悉感与滋养} & 停止探索新世界，停止焦虑未来 & 吃熟悉的舒适食物，温水澡，简化待办，回安全屋\tabularnewline
Ne & ISTJ, ISFJ & \textbf{现实检验锚点} & 停止灾难化想象，停止过度推演 & 罗列已知事实，找人分担，做微小但确定的改变\tabularnewline
Ni & ESTP, ESFP & \textbf{放过明天} & 停止当下反应，停止试图解决未来 & 推迟所有决策，借用他人愿景，纯粹的娱乐转移\tabularnewline
\midrule
\multicolumn{5}{c}{判断解药：回归秩序与人性 (Judgement)}\tabularnewline
\midrule
Te & INFP, ISFP & \textbf{小赢策略} & 停止宏大目标，停止自我攻击 & 整理一个抽屉，划掉简单任务，整理房间，承认无力\tabularnewline
Ti & ENFJ, ESFJ & \textbf{逻辑隔离} & 停止情绪读心，停止人际纠结 & 独处充电，玩解谜/数字游戏，做不涉及人的智力活动\tabularnewline
Fe & INTP, ISTP & \textbf{无压陪伴} & 停止彻底隔离，停止逻辑辩论 & 并在但各玩各的(Parallel Play)，书面沟通，身体接触\tabularnewline
Fi & ENTJ, ESTJ & \textbf{允许脆弱} & 停止压抑分析，停止强撑坚强 & 私密环境哭泣宣泄，暂停工作，寻找安全盟友倾诉\tabularnewline
\bottomrule
\end{tabular}}
\end{table}


\begin{table}[h!]
\caption[认知功能的5*反向作用(Opposing/Nemesis)表现：防御与质疑]{认知功能的5*反向作用(Opposing/Nemesis)表现：防御与质疑。
这是进入潜意识的第一道防线。当我们的主导功能(Hero)受到挑战时，第5功能会跳出来“抬杠”。表现形式：固执、怀疑论、被动攻击，或者对他人的该功能表现感到莫名的厌烦。}
    \label{tab:cog-function-opposing}
\centering
\resizebox{\textwidth}{!}{
\begin{tabular}{ccccc}
\toprule
反向功能 & 涉及类型 & 核心潜台词 & 防御表现 (Stubbornness) & 投射怀疑 (Paranoia)\tabularnewline
\midrule
\multicolumn{5}{c}{感知型反向：拒绝接受信息}\tabularnewline
\midrule
Ne & \cellcolor{intjpurple}\color{white}INTJ, INFJ & “别扯那些没用的” & \textbf{厌恶变数}：固执地坚持原计划，拒绝听取其他可能性，认为那是干扰 & \textbf{怀疑意图}：总觉得别人背信弃义，或者怀疑别人的忠诚度\tabularnewline
Ni & ENTP, ENFP & “你想控制我？” & \textbf{抗拒定论}：讨厌被“贴标签”或被限制未来，故意反其道而行之 & \textbf{怀疑操控}：认为别人的愿景是洗脑，怀疑背后有阴谋\tabularnewline
Se & ISTJ, ISFJ & “我不想看/听” & \textbf{感官屏蔽}：对外界的新体验表现出排斥，固守熟悉的物理环境 & \textbf{怀疑现实}：认为别人的衣着/表现是肤浅的作秀\tabularnewline
Si & ESTP, ESFP & “别拿过去压我” & \textbf{轻视经验}：极度讨厌“老规矩”或繁琐的流程，故意打破常规 & \textbf{怀疑动机}：认为别人拿过去的数据是在限制自己的自由\tabularnewline
\midrule
\multicolumn{5}{c}{判断型反向：拒绝接受逻辑/价值}\tabularnewline
\midrule
Ti & ENTJ, ESTJ & “这没意义” & \textbf{傲慢驳斥}：若是不能马上应用，就拒绝深入思考逻辑细节，认为那是浪费时间 & \textbf{怀疑智商}：认为别人的精细逻辑是“书呆子气”，试图用权威压倒道理\tabularnewline
Te & INTP, ISTP & “你在教我做事？” & \textbf{抗拒指挥}：极度厌恶被别人安排或催促，故意拖延以示反抗 & \textbf{怀疑能力}：认为别人的证书/头衔都是虚的，质疑权威的有效性\tabularnewline
Fi & ENFJ, ESFJ & “你太自私了” & \textbf{否定个体}：为了群体和谐而压抑个人感受，也看不惯别人的特立独行 & \textbf{怀疑人品}：认为那些坚持自我价值的人是“破坏团结”的坏人\tabularnewline
Fe & INFP, ISFP & “别以此绑架我” & \textbf{拒绝从众}：极度反感社交压力和“必须做的事”，故意表现得冷漠 & \textbf{怀疑虚伪}：认为别人的客套和热情都是有目的的操纵\tabularnewline
\bottomrule
\end{tabular}}
\end{table}


\begin{table}
\caption[认知功能的6*批判家长(Critical Parents)表现]{认知功能的6*批判家长(Critical Parents)表现: 如何打压别人和自己}
    \label{tab:cog-function-critical}
    \resizebox{\textwidth}{!}{
\begin{tabular}{ccccc}
\toprule
批判功能 & 涉及类型 & 核心台词 & 对外打压 & 对内自毁\tabularnewline
\midrule
\multicolumn{5}{c}{理性审判组：Judgement}\tabularnewline
\midrule
Ti Critical 批判逻辑 & INTJ, ISTJ & 你真蠢 & 智力羞辱，冷漠指出他人逻辑漏洞 & 自我怀疑，质疑自己不够聪明\tabularnewline
Fi Critical 批判道德 & INFJ, ISFJ & 你真坏 & 道德审判，指责自私虚伪冷血人品 & 自我羞耻，觉得自己是伪君子\tabularnewline
Te Critical 批判能力 & ENTP, ESTP & 你没用 & 能力否定，嘲笑笨手笨脚一事无成 & 成就焦虑，成功仍觉一无是处\tabularnewline
Fe Critical 批判社交 & ENFP, ESFP & 没人理你 & 社交孤立，攻击对方情商或者礼仪 & 归属崩塌，总觉自己格格不入\tabularnewline
\midrule
\multicolumn{5}{c}{感知审判组：Perception}\tabularnewline
\midrule
Ni Critical 批判远见 & INTP, INFP & 没希望 & 未来诅咒，否定愿景断言注定失败 & 虚无主义，觉得做什么都没意义\tabularnewline
Si Critical 批判经验 & ISTP, ISFP & 你真弱 & 翻旧帐，指责不可靠不长记性/否定毅力 & 创伤反刍，不断回忆过去的痛苦\tabularnewline
Ne Critical 批判意图 & ENTJ, ENFJ & 你想害我 & 动机阴谋论，怀疑忠诚指责对方鲁莽 & 选择障碍，看到无数会感到受困\tabularnewline
Se Critical 批判现实 & ESTJ, ESFJ & 真难看 & 外貌/品味攻击，苛刻挑剔对方衣着审美 & 容貌/现实焦虑，觉得自己丑陋笨拙\tabularnewline
\bottomrule
\end{tabular}}
\end{table}





\begin{table}[h!]
\caption[认知功能的7*盲点(Trickster)表现：无意识的黑洞]{认知功能的7*盲点(Trickster)表现：无意识的黑洞。
我们不仅不擅长此功能，而且往往会通过贬低它来保护自尊。表现形式：根本看不见其存在的必要性，将其视为多余、虚伪或荒谬的干扰。}
    \label{tab:cog-function-trickster}
\centering
\resizebox{\textwidth}{!}{
\begin{tabular}{cccc}
\toprule
盲点功能 & 涉及类型 & 盲区心态/潜台词 (The Dismissal) & 盲点具象表现\tabularnewline
\midrule
\multicolumn{4}{c}{对人际与道德的盲视}\tabularnewline
\midrule
Fe & \cellcolor{intjpurple}\color{white}INTJ, ISTJ & “虚伪的客套有什么用？讲重点。” & 视社交礼仪为操纵，善意说真话却成冒犯，完全读不懂空气\tabularnewline
Fi & ENTP, ESTP & “这也太矫情了，客观一点行不行？” & 情感逻辑化，不知道自己真实感受，无意间把别人的价值观踩得粉碎\tabularnewline
\midrule
\multicolumn{4}{c}{对现实与身体的盲视}\tabularnewline
\midrule
Se & INTP, INFP & “这重要吗？我没看见/没听见。” & 物理笨拙(平地摔)，极其忽视环境细节，活在脑洞里完全缺乏行动力\tabularnewline
Si & ENTJ, ENFJ & “别提过去，我只要未来！” & 严重透支身体而不自知，极度厌恶怀旧和重复，记不住生活琐事\tabularnewline
\midrule
\multicolumn{4}{c}{对联想与愿景的盲视}\tabularnewline
\midrule
Ne & ISTP, ISFP & “你到底想说什么？别扯那些没用的。” & 极度字面化，理解不了暗示和隐喻，讨厌猜测，一条路走到黑\tabularnewline
Ni & ESTJ, ESFJ & “别想太多，先做再说，这不切实际。” & 短期主义，只要确定的流程而不要愿景，瞎操心，对未知充满恐惧\tabularnewline
\midrule
\multicolumn{4}{c}{对逻辑与效率的盲视}\tabularnewline
\midrule
Te & INFJ, ISFJ & “你这人怎么这么冷血/功利？” & 无法理清外部资源，难以做客观决策，为了维护人际或信仰而抗拒数据\tabularnewline
Ti & ENFP, ESFP & “太较真了没意思，差不多就行了。” & 逻辑疏松，讨厌定义和辩论，经常自我矛盾却不自知，只关注“好不好”\tabularnewline
\bottomrule
\end{tabular}}
\end{table}


\begin{table}

\caption[认知功能的8*恶魔(Demon)表现]{认知功能的8*恶魔(Demon)表现： 彻底崩溃时的毁灭打击。这是人格结构中最深层，最黑暗，也是最原始的部分。触发条件：自尊彻底崩溃，核心界限被严重侵犯，或者感受到极度绝望。表现形式：毁灭一切，不在乎输赢，只想让一切归零
-- “既然你毁了我的世界，那我就把整个宇宙都炸了陪葬”。转化潜力：如果能整合，它也能带来涅槃重生的力量(The Daemonic
vs. The Demonic)。}
    \label{tab:cog-function-demon}
    \centering
\resizebox{\textwidth}{!}{
\begin{tabular}{cccc}
\toprule
恶魔功能 & 涉及类型 & 核心心态 & 毁灭性表现(The Destruction)\tabularnewline
\midrule
\multicolumn{4}{c}{感知型恶魔}\tabularnewline
\midrule
Si Demon & \cellcolor{intjpurple}\color{white}INTJ, INFJ & 彻底抹除 | Erasure & 记忆清洗与断连(The Door Slam)，肉体自毁倾向\tabularnewline
Se Demon & ENTP, ENFP & 现实崩塌 | Burn it down & 物理毁灭与虚无，极度慵懒或狂躁\tabularnewline
Ni Demon & ISTJ, ISFJ & 永恒绝望 | Hopelessness & 恐怖末日预言家，否定别人的梦想\tabularnewline
Ne Demon & ISTP, ISFP & 恶意扭曲 | Gaslighting & 被害妄想与构陷，混乱制造者\tabularnewline
\midrule
\multicolumn{4}{c}{判断型恶魔}\tabularnewline
\midrule
Ti Demon & INFP, ISFP & 诛心真话 | Cold Truth & 最恶毒的逻辑穿刺，摧毁对方自尊\tabularnewline
Te Demon & INTP, ISTP & 系统毁灭 | System Crash & 冷酷的暴君，破坏性执行，无视人性\tabularnewline
Fi Demon & ENTP, ESTP & 道德虚无 | Nihilism & 反社会倾向，情感报复\tabularnewline
Fe Demon & ENTJ, ESTJ & 名誉屠杀 | Character Assassination & 情感操纵与社会死亡，情感爆发\tabularnewline
\bottomrule
\end{tabular}}
\end{table}




\begin{table}

\caption[认知功能的8*恶魔(Demon)整合]{认知功能的8*恶魔(Demon)整合：涅槃重生后的救赎之力(The Daemonic)。当恶魔功能被意识整合，它不再是毁灭者，而变成了通往完整人格的守门人。这种力量通常带有一种“非个人”的、超越小我的神性或宿命感。表现形式：大智若愚、绝地反击、或是对他人的深刻救赎。}
    \label{tab:cog-function-demon-redemption}
    \centering

\resizebox{\textwidth}{!}{
\begin{tabular}{cccc}
\toprule
恶魔功能 & 涉及类型 & 转化核心 & 整合后表现(The Redemption)\tabularnewline
\midrule
\multicolumn{4}{c}{感知型恶魔}\tabularnewline
\midrule
Si Demon & \cellcolor{intjpurple}\color{white}INTJ, INFJ & 历史和解 | Reconciliation & 成为传统的守护者，从过往创伤中提炼史诗般的意义，极致的身体掌控\tabularnewline
Se Demon & ENTP, ENFP & 绝对显化 | Manifestation & 将抽象疯狂的灵感在现实中完美落地，极具感染力的舞台表现，活在当下\tabularnewline
Ni Demon & ISTJ, ISFJ & 先知智慧 | Visionary & 突破经验主义的桎梏，在危机时刻涌现出精准的未来预判，指引方向\tabularnewline
Ne Demon & ISTP, ISFP & 万物互联 | Interconnection & 极其通透的开放心态，理解复杂的混沌系统，成为打破僵局的创新者\tabularnewline
\midrule
\multicolumn{4}{c}{判断型恶魔}\tabularnewline
\midrule
Ti Demon & INFP, ISFP & 结构真理 | Structural Truth & 能够用最客观的逻辑捍卫价值观，不带情绪地通过理性解决最棘手的问题\tabularnewline
Te Demon & INTP, ISTP & 建设秩序 | Constructive Order & 极具执行力的系统构建者，在混乱中建立坚不可摧的帝国，为了大局的高效\tabularnewline
Fi Demon & ENTP, ESTP & 慈悲圣徒 | Universal Love & 浪子回头般的真诚，对他人的苦难产生深刻共情，为了信仰不惜一切\tabularnewline
Fe Demon & ENTJ, ESTJ & 仁慈领袖 | Benevolent Ruler & 放下权力的傲慢，展现出真正的人性光辉，通过情感连接凝聚人心\tabularnewline
\bottomrule
\end{tabular}}
\end{table}










\begin{table}[h!]
\caption[认知功能的防御机制/长期内耗循环(Loop)]{认知功能的防御机制/长期内耗循环(Loop)。不同于瞬间压力的Grip，Loop是一种长期的、病态的舒适区，个体切断了辅助功能(第二功能)的平衡作用，只在第一和第三功能之间死循环。}
    \label{tab:cog-function-loop}
    \centering
    \resizebox{\textwidth}{!}{
\begin{tabular}{cccc}
\toprule
Loop组合 & 状态名称 & 涉及类型 & 典型表现 (Manifestation)\tabularnewline
\midrule
\multicolumn{4}{c}{内向Loop：切断外部互动 $\rightarrow$ 陷入大脑死循环 $\rightarrow$ 拒绝现实检验}\tabularnewline
\midrule
\multirow{2}{*}{Ni - Fi} & \multirow{2}{*}{偏执内耗 (Paranoid)} & \cellcolor{intjpurple}\color{white}INTJ & 沉溺消极愿景，觉得自己是受害者，全世界都在针对自己，无行动验证\tabularnewline
 & & ISFP & 陷入深刻的负面自我定义，确信无人能理解自己，自我孤立\tabularnewline
\midrule
\multirow{2}{*}{Ti - Si} & \multirow{2}{*}{固步自封 (Stuck)} & INTP & 反复分析过去错误，抗拒新信息，甚至忽略饮食卫生等基本生存需求\tabularnewline
 & & ISFJ & 翻旧账大师，逻辑证明自己注定失败，拒绝任何新尝试\tabularnewline
\midrule
\multirow{2}{*}{Ni - Ti} & \multirow{2}{*}{虚无主义 (Nihilistic)} & INFJ & 极度冷漠愤世嫉俗，切断同情心，构建宏大理论证明“人类没救了”\tabularnewline
 & & ISTP & 对未来产生阴谋论，逻辑自洽地认为别人甚至宇宙都在算计自己\tabularnewline
\midrule
\multirow{2}{*}{Fi - Si} & \multirow{2}{*}{受害循环 (Resentful)} & INFP & 强迫性重温过去创伤，认为未来只会重复痛苦，丧失所有希望\tabularnewline
 & & ISTJ & 变得情绪化且敏感，过度解读过去的细节，赋予其沉重的负面色彩\tabularnewline
\midrule
\multicolumn{4}{c}{外向Loop：切断内心联系 $\rightarrow$ 疯狂外部活动 $\rightarrow$ 逃避深度思考}\tabularnewline
\midrule
\multirow{2}{*}{Te - Se} & \multirow{2}{*}{鲁莽暴君 (Bulldozer)} & ENTJ & 为行动而行动，停不下来，通过高强度的感官/工作麻痹思考，独断专行\tabularnewline
 & & ESFP & 极度强势急功近利，不仅自己享乐，还强迫别人服从自己的“效率”\tabularnewline
\midrule
\multirow{2}{*}{Ne - Fe} & \multirow{2}{*}{表演人格 (Manipulative)} & ENTP & 极度渴望关注，失去逻辑锚点，为获得反应故意挑衅或成为空洞梗王\tabularnewline
 & & ESFJ & 脑洞大开地焦虑他人看法，变得八卦、戏剧化，为表面和谐撒谎\tabularnewline
\midrule
\multirow{2}{*}{Fe - Se} & \multirow{2}{*}{浅薄迎合 (Superficial)} & ENFJ & 沉迷当下感官与评价，冲动消费，过度社交，忘记长远目标\tabularnewline
 & & ESTP & 虚伪的社交变色龙，为了群体认可过度表现情感，在不同圈子演戏\tabularnewline
\midrule
\multirow{2}{*}{Ne - Te} & \multirow{2}{*}{焦虑忙碌 (Burnout)} & ENFP & 开启一百个计划却无一完成，强迫性快速执行，无法容忍情感反思\tabularnewline
 & & ESTJ & 异想天开却独断专行，提出不切实际的新规并强迫大家执行，抛弃经验\tabularnewline
\bottomrule
\end{tabular}}
\end{table}



\begin{table}[h!]
\caption[长期内耗循环(Loop)的解药]{长期内耗循环(Loop)的解药。核心逻辑：Loop 是因为忽略了辅助功能（家长），导致第一与第三功能恶性循环。解药是\textbf{强制调用第二功能}。}
    \label{tab:cog-function-loop-remedy}
    \centering
    \resizebox{\textwidth}{!}{
\begin{tabular}{cccc}
\toprule
解药 (辅助功能) & 涉及 Loop & 核心逻辑 (Logic) & 行动指南 (Actionable Steps)\tabularnewline
\midrule
\multicolumn{4}{c}{内向 Loop 解药：向外连接，破除自我封闭 (Externalize)}\tabularnewline
\midrule
\multirow{2}{*}{Te (外倾思考)} & \cellcolor{intjpurple}\color{white}Ni-Fi (INTJ) & \textbf{产出与秩序} & \multirow{2}{*}{列出清单，整理物理环境，追求“小赢”，将担忧转化为待办}\tabularnewline
 & Si-Fi (ISTJ) & 停止感受，开始执行 & \tabularnewline
\midrule
\multirow{2}{*}{Fe (外倾情感)} & Ni-Ti (INFJ) & \textbf{表达与连接} & \multirow{2}{*}{找人倾诉，进行辅导/帮助他人，寻求外部的情感确认}\tabularnewline
 & Si-Ti (ISFJ) & 停止自剖，关注他人 & \tabularnewline
\midrule
\multirow{2}{*}{Se (外倾感觉)} & Ti-Ni (ISTP) & \textbf{行动与感官} & \multirow{2}{*}{高强度运动，手工制造，大扫除，接触自然，强迫回到当下}\tabularnewline
 & Fi-Ni (ISFP) & 停止空想，物理接触 & \tabularnewline
\midrule
\multirow{2}{*}{Ne (外倾直觉)} & Ti-Si (INTP) & \textbf{探索与发散} & \multirow{2}{*}{物理换环境，摄入随机信息(书/剧)，头脑风暴，尝试新路线}\tabularnewline
 & Fi-Si (INFP) & 停止反刍，寻找新意 & \tabularnewline
\midrule
\multicolumn{4}{c}{外向 Loop 解药：内向沉淀，停下来检查内心 (Internalize)}\tabularnewline
\midrule
\multirow{2}{*}{Ni (内倾直觉)} & Te-Se (ENTJ) & \textbf{愿景与暂停} & \multirow{2}{*}{强制独处，断网冥想，推演长远后果，追问做事的根本意义}\tabularnewline
 & Fe-Se (ENFJ) & 停止忙碌，寻找本质 & \tabularnewline
\midrule
\multirow{2}{*}{Si (内倾感觉)} & Te-Ne (ESTJ) & \textbf{复盘与稳健} & \multirow{2}{*}{查阅历史先例，关注身体感受，做可行性检查，休息与维护}\tabularnewline
 & Fe-Ne (ESFJ) & 停止折腾，回归经验 & \tabularnewline
\midrule
\multirow{2}{*}{Fi (内倾情感)} & Ne-Te (ENFP) & \textbf{价值与本心} & \multirow{2}{*}{日记反思，道德审查，学会说“不”，确认是否符合真实喜好}\tabularnewline
 & Se-Te (ESFP) & 停止取悦，自问本心 & \tabularnewline
\midrule
\multirow{2}{*}{Ti (内倾思考)} & Ne-Fe (ENTP) & \textbf{逻辑与真理} & \multirow{2}{*}{独立逻辑分析，精确定义概念，消除模糊，不论人情只论对错}\tabularnewline
 & Se-Fe (ESTP) & 停止表演，回归事实 & \tabularnewline
\bottomrule
\end{tabular}}
\end{table}





\begin{table}[h!]
\caption[状态(State) vs. 特质(Trait)：预期管理表]{状态(State) vs. 特质(Trait)：预期管理表。
绝大多数人(80\%+)长期处于轻微 Loop 或 Grip 中，这是现代生活的常态。区分的关键在于“可逆性”与“自我认知”。}
\centering
\resizebox{\textwidth}{!}{
\begin{tabular}{cccc}
\toprule
维度 & 暂时状态 (State) & 固化特质 (Trait) & INTJ 的预期管理策略\tabularnewline
\midrule
\textbf{定义} & 应激反应 (Stress Response) & 人格结构 (Structure) & 哪怕是圣人也会有 Loop 时刻\tabularnewline
\midrule
\textbf{自我认知} & \textbf{自我失调 (Ego-dystonic)} & \textbf{自我协调 (Ego-syntonic)} & 观察他事后是否会道歉或反思\tabularnewline
& 觉得“我不该这样，这不像我” & 觉得“我没错，是世界的错” & 只有前者值得深入交往\tabularnewline
\midrule
\textbf{触发机制} & 有明确的外部压力源 & 无需触发，是默认行为模式 & 寻找触发点，若无因却疯癫，则远离\tabularnewline
& (如：失业、失恋、过劳) & (如：自恋、边缘、反社会) & \tabularnewline
\midrule
\textbf{表现形式} & \textbf{Grip/Loop} & \textbf{人格障碍 (PD)} & 区分“受伤的野兽”与“原本的恶魔”\tabularnewline
& 情绪过后会回归 Hero/Parent & 阴影功能喧宾夺主，成为常态 & 前者需要空间，后者需要断联\tabularnewline
\midrule
\textbf{恢复弹性} & \textbf{高弹性} & \textbf{无弹性/僵化} & 不要试图“拯救”特质固化的人\tabularnewline
& 给点解药(如休息)就能弹回 & 拒绝改变，拒绝辅助功能介入 & 那是心理医生的工作，不是朋友的\tabularnewline
\bottomrule
\end{tabular}}
\end{table}


\begin{table}[h!]
\caption[全谱系战略应对表：针对 8 种内耗循环(Loop)的社交博弈策略]{全谱系战略应对表：针对 8 种内耗循环(Loop)的社交博弈策略。
核心心法：诊断出 Loop 后，切忌在对方的逻辑里打转。内向 Loop 需要“外化引导”，外向 Loop 需要“边界阻断”。}
\centering
\resizebox{\textwidth}{!}{
\begin{tabular}{ccccc}
\toprule
Loop 类型 & 涉及人格 & 对方心态 (The Mindset) & 你的错误反应 (Don't) & 战略反应 (Do: 降维打击)\tabularnewline
\midrule
\multicolumn{5}{c}{外向 Loop：躁动、表演、强迫执行 (对策：设立边界，强制降速)}\tabularnewline
\midrule
\textbf{Te-Se} & ENTJ & “我要赢，立刻，马上” & “你太功利了，休息一下” & \textbf{顺势借力 (Aikido)}：\tabularnewline
鲁莽暴君 & ESFP & (停不下来的推土机) & (试图让他慢下来) & 承认他的效率，给他一个极难的具体任务消耗其过剩精力\tabularnewline
\midrule
\textbf{Fe-Se} & ENFJ & “快看我，夸我，爱我” & “你太虚伪了，是个戏精” & \textbf{情绪喂养/隔离 (Grey Rock)}：\tabularnewline
浅薄迎合 & ESTP & (活在他人眼里的空心人) & (当面揭穿或冷脸) & 像喂金鱼一样给点低成本赞美(Feed Fe)，然后迅速物理撤离\tabularnewline
\midrule
\textbf{Ne-Fe} & ENTP & “我要搞个大新闻/看反应” & “你逻辑不通/我很生气” & \textbf{逻辑锚点 (Logic Anchor)}：\tabularnewline
表演人格 & ESFJ & (混乱善良/邪恶的乐子人) & (给他提供情绪反馈) & 不给情绪反应，冷漠地追问定义：“你刚才那个词具体指什么？”\tabularnewline
\midrule
\textbf{Ne-Te} & ENFP & “我有100个计划你必须配合” & “你这根本不落地/不现实” & \textbf{价值排序 (Prioritize)}：\tabularnewline
焦虑忙碌 & ESTJ & (暴躁的无头苍蝇) & (试图用经验Si阻拦) & 不否认计划，但问核心：“哪个最重要？我们先只做这一个。”\tabularnewline
\midrule
\multicolumn{5}{c}{内向 Loop：封闭、偏执、自我反刍 (对策：打破真空，物理介入)}\tabularnewline
\midrule
\textbf{Ni-Fi} & \cellcolor{intjpurple}\color{white}INTJ & “众人皆醉，世界没救了” & “别想太多，出来玩” & \textbf{任务外化 (Externalize)}：\tabularnewline
偏执内耗 & ISFP & (深陷消极愿景的受害者) & (试图用Se强行拉扯) & 别辩论观点。把话题变成物体：“这事依然无解，但能帮我递下杯子吗？”\tabularnewline
\midrule
\textbf{Ti-Si} & INTP & “过去失败证明未来不可行” & “你要勇敢尝试/别钻牛角尖” & \textbf{输入变量 (New Input)}：\tabularnewline
固步自封 & ISFJ & (翻旧账的逻辑死循环) & (攻击他的逻辑闭环) & “你分析得对。但如果加入这个新参数(Ne)，结果会变吗？”\tabularnewline
\midrule
\textbf{Ni-Ti} & INFJ & “我在揭露宇宙/人性的黑暗” & “你太消极/阴谋论了” & \textbf{现实检验 (Reality Check)}：\tabularnewline
虚无主义 & ISTP & (愤世嫉俗的旁观者) & (试图用正能量说教) & “这理论很完美。但它能解释为什么那个人刚刚对你笑了吗？”(引导看Se)\tabularnewline
\midrule
\textbf{Fi-Si} & INFP & “我永远走不出这个创伤” & “这点小事至于吗/向前看” & \textbf{允许哀悼 (Validation)}：\tabularnewline
受害循环 & ISTJ & (沉溺于过去的痛苦细节) & (试图提供Te解决方案) & 不要给建议。只要点头确认：“那确实很痛。”(等到情绪耗尽再谈Te)\tabularnewline
\bottomrule
\end{tabular}}
\end{table}


\begin{table}[h!]
\caption[全谱系社会控制矩阵：资本主义如何无死角收割 8 种认知循环(Loop)]{全谱系社会控制矩阵：资本主义如何无死角收割 8 种认知循环(Loop)。
核心逻辑：Loop 是个体在逃避痛苦时的防御机制，而市场经济将每一种“防御”都包装成了“消费需求”。}
\centering
\resizebox{\textwidth}{!}{
\begin{tabular}{ccccc}
\toprule
Loop 类型 & 涉及人格 & 资本主义的话术 (The Pitch) & 收割方式/产品 (The Product) & 典型受害者心态\tabularnewline
\midrule
\multicolumn{5}{c}{外向 Loop 收割：利用欲望、焦虑与虚荣}\tabularnewline
\midrule
\textbf{Te-Se} & ENTJ & “狼性文化” & \textbf{效率奢侈品 \& 成功学} & “我必须更强、更快、更富。”\tabularnewline
鲁莽暴君 & ESFP & “Work Hard, Play Hard” & 豪车名表、能量饮料、高端商务舱、速成课 & (停不下来的生产/消费机器)\tabularnewline
\midrule
\textbf{Fe-Se} & ENFJ & “精致生活” & \textbf{颜值经济 \& 社交货币} & “如果不美/不合群，我就没价值。”\tabularnewline
浅薄迎合 & ESTP & “对自己好一点” & 医美贷、网红打卡、快时尚、社交会员 & (活给别人看的流量奴隶)\tabularnewline
\midrule
\textbf{Ne-Fe} & ENTP & “错过即过错 (FOMO)” & \textbf{注意力收割 \& 炒作经济} & “大家都在玩这个，我不能落下。”\tabularnewline
表演人格 & ESFJ & “跟上潮流/热梗” & 短视频算法、Meme币、盲盒、NFT、热搜 & (被算法随意调度的跟风者)\tabularnewline
\midrule
\textbf{Ne-Te} & ENFP & “副业刚需/斜杠青年” & \textbf{焦虑贩卖 \& 生产力工具} & “我还要做更多，睡觉就是浪费。”\tabularnewline
焦虑忙碌 & ESTJ & “多任务并行/优化人生” & 知识付费、Notion模板、理财课、抢票加速包 & (被“机会”累死的无头苍蝇)\tabularnewline
\midrule
\multicolumn{5}{c}{内向 Loop 收割：利用恐惧、孤独与傲慢}\tabularnewline
\midrule
\textbf{Ti-Si} & INTP & “未雨绸缪” & \textbf{恐慌经济 \& 安全感} & “世界太危险，必须囤积一切。”\tabularnewline
固步自封 & ISFJ & “规避风险/养老焦虑” & 保险、黄金、防脱发、囤积癖(Storage)、硬核养生 & (为“可能发生的灾难”买单)\tabularnewline
\midrule
\textbf{Fi-Si} & INFP & “治愈你的内在小孩” & \textbf{怀旧经济 \& 疗愈消费} & “现实太痛了，我要躲回过去。”\tabularnewline
受害循环 & ISTJ & “还是过去的味道好” & 复古潮牌、解压玩具、甚至被滥用的心理咨询/灵修 & (沉溺于创伤抚慰的巨婴)\tabularnewline
\midrule
\textbf{Ni-Fi} & \cellcolor{intjpurple}\color{white}INTJ & “独立思考/品味隔离” & \textbf{鄙视链顶端 \& 智商税} & “只有我懂它的价值，你们都是俗人。”\tabularnewline
偏执内耗 & ISFP & “为信仰/本质充值” & 极简主义天价单品、绝版收藏、Hi-Fi音响、徕卡 & (为“独特性”支付溢价的孤岛)\tabularnewline
\midrule
\textbf{Ni-Ti} & INFJ & “看清世界的真相” & \textbf{阴谋论 \& 替代性知识} & “主流都是谎言，我要寻找解药。”\tabularnewline
虚无主义 & ISTP & “跳出矩阵 (Red Pill)” & 区块链信仰、甚至邪教课程、生存狂装备(Prepper) & (试图智力碾压系统的边缘人)\tabularnewline
\bottomrule
\end{tabular}}
\end{table}


\begin{table}[h!]
\caption[INTJ 的四重人格状态切换表：你体内的隐藏角色]{INTJ 的四重人格状态切换表：你体内的隐藏角色。
核心原理：当主我(Ego)疲惫、自信爆棚或极度受挫时，我们会切换到其他象限接管身体。
\newline
\textbf{阴影吞噬机制详解} \\
\textbf{1. 解释}：阴影吞噬不仅仅是压力的表现，而是Ego（自我）的完全崩溃。此时，意识的主导权移交给了第5-8功能，个体表现出与平时截然相反的“反面人格”，通常带有极强的破坏性、偏执和幼稚感。 \\
\textbf{2. 触发原理（Trigger）}：长期且巨大的心理压力导致第4功能（劣势功能）防御失效 $\rightarrow$ 心理大坝决堤 $\rightarrow$ 潜意识阴影全线接管。 \\
\textbf{3. 应对方式}：\textit{切勿试图讲理或对抗}。处于此状态的人逻辑是扭曲的（第6功能批判家）且甚至由于现实感丧失（第8功能魔鬼）而不可理喻。唯一解法是：\textbf{物理隔离 + 极端的感官休息（睡眠/独处）}，等待能量燃尽后的重启。}

\centering
\resizebox{\textwidth}{!}{
\begin{tabular}{ccccc}
\toprule
状态名称 & 对应人格 & 触发条件 & 典型表现 (Mode) & INTJ 的体验\tabularnewline
\midrule
\textbf{本我 (Ego)} & \textbf{INTJ} & \textbf{默认状态} & \textbf{战略家 (The Strategist)} & 冷静、规划、内省、效率、低能量消耗\tabularnewline
\midrule
\textbf{潜意识 (Subconscious)} & \textbf{ESFP} & \textbf{极度自信/放松} & \textbf{表演者 (The Performer)} & \textbf{这解释了你的“高能量”时刻}！\tabularnewline
(Hero的反面) & (Aspiring) & 或需要展示魅力时 & (Se-Fi 主导) & 当你拿起相机(Se)进入心流，不再焦虑未来(Ni)时，\tabularnewline
& & & & 你会变得幽默、活在当下、极具感染力。\tabularnewline
\midrule
\textbf{无意识 (Unconscious)} & \textbf{ENTP} & \textbf{防御/被质疑} & \textbf{辩论家 (The Debater)} & 变得多疑、爱抬杠、为了赢而赢。\tabularnewline
(Shadow 阴影) & & 智力受到挑战时 & (Ne-Ti 主导) & Ne Nemesis 让你怀疑一切，Ti Critic 让你刻薄地攻击别人的逻辑漏洞。\tabularnewline
\midrule
\textbf{超我 (Superego)} & \textbf{ISFJ} & \textbf{极度绝望/毁灭} & \textbf{复仇天使 (The Punisher)} & \textbf{地狱模式}。彻底放弃逻辑和愿景，\tabularnewline
(Demon 恶魔) & & 核心价值被摧毁时 & (Si-Fe 主导) & 变得极度计较细节(Si Demon)，想要毁灭世界或彻底自我牺牲(Fe)。\tabularnewline
\bottomrule
\end{tabular}}
\end{table}


\begin{table}[h!]
\caption[社交狩猎矩阵：PUA与精神控制如何利用认知功能漏洞]{社交狩猎矩阵：PUA与精神控制是如何利用认知功能漏洞的。
核心原理：操纵者通过攻击受害者的“劣势功能”引发焦虑，再通过扮演“完美拯救者”进行收割。不同类型的痛点（软肋）完全不同。}
\centering
\resizebox{\textwidth}{!}{
\begin{tabular}{ccccc}
\toprule
受害目标 & 软肋功能 & 典型操纵手段 (The Tactic) & 攻击话术/心理暗示 & 为什么会生效 (Mechanism)\tabularnewline
\midrule
\multicolumn{5}{c}{针对 NT 分析师：利用“傲慢”与“社交无能”}\tabularnewline
\midrule
\rowcolor{intjpurple!20} \textbf{INTJ} & \textbf{Se 劣势} & \textbf{Negging (打压/否定)} & “你确实很聪明，但有点无趣/死板。” & \textbf{攻击 Se 表演欲}：INTJ 潜意识渴望Se的潇洒，\tabularnewline
\rowcolor{intjpurple!20} ENTJ & \textbf{Fi 劣势} & \& \textbf{冷读术} & “我知道你外表冷漠，其实内心渴望狂野。” & 一旦被定义为“无趣”，就会拼命通过顺从对方来证明自己。\tabularnewline
\midrule
\textbf{INTP} & \textbf{Fe 劣势} & \textbf{Frozen (冷冻/撤回)} & “既然你这么想，那随你便吧。(沉默)” & \textbf{制造 Fe 恐慌}：利用 TiUser 对人际氛围的不知所措，\tabularnewline
ENTP & \textbf{Si 劣势} & \& \textbf{智力隔离} & “这点常识都不懂？我没法跟你解释。” & 突然撤回关注，迫使对方为了恢复“和谐”而放弃逻辑底线。\tabularnewline
\midrule
\multicolumn{5}{c}{针对 NF 外交官：利用“圣母心”与“自我怀疑”}\tabularnewline
\midrule
\textbf{INFJ} & \textbf{Se 劣势} & \textbf{Gaslighting (煤气灯)} & “你记错了，根本没发生过，你太敏感了。” & \textbf{瓦解 Se 现实感}：利用 NF 对物理现实的不确定，\tabularnewline
ENFJ & \textbf{Ti 劣势} & \& \textbf{受害者扮演} & “都是因为爱你我才疯的，你却在分析我？” & 篡改记忆。同时利用 Fe 的同情心，让受害者对施暴者产生负疚感。\tabularnewline
\midrule
\textbf{INFP} & \textbf{Te 劣势} & \textbf{Dependence (习得性无助)} & “没有我，你连生活都不能自理。” & \textbf{打击 Te 能力感}：不断强调受害者在现实层面的无能，\tabularnewline
ENFP & \textbf{Si 劣势} & \& \textbf{未来画饼} & “忍过这一阵，我们的未来会很美。” & 使其确信离开操纵者就无法生存。利用 Ne 给人造一个虚假的未来。\tabularnewline
\midrule
\multicolumn{5}{c}{针对 SJ 守护者：利用“责任感”与“恐惧变化”}\tabularnewline
\midrule
\textbf{ISTJ} & \textbf{Ne 劣势} & \textbf{Fear-Mongering (贩卖恐惧)} & “如果你离开我，你会孤独终老/一无所有。” & \textbf{引爆 Ne 灾难想象}：利用 SJ 对未知的极度恐惧，\tabularnewline
ESTJ & \textbf{Fi 劣势} & \& \textbf{道德绑架} & “你怎么能这么自私？我不符合规矩吗？” & 让他们不敢踏出舒适区半步。利用 Fi 的道德枷锁使其服从。\tabularnewline
\midrule
\textbf{ISFJ} & \textbf{Ne 劣势} & \textbf{Love Bombing (爱意轰炸)} & (初期)“你是完美的女神” $\to$ (后期)撤回 & \textbf{多巴胺戒断}：先满足 Fe/Si 的被需要感，建立依赖，\tabularnewline
ESFJ & \textbf{Ti 劣势} & \& \textbf{三角测量} & “前任比你懂事多了/别人都觉得你不对。” & 再引入竞争者制造焦虑，让其陷入“为了赢回爱”的内卷中。\tabularnewline
\midrule
\multicolumn{5}{c}{针对 SP 探索者：利用“自由”与“智力自卑”}\tabularnewline
\midrule
\textbf{ISTP} & \textbf{Fe 劣势} & \textbf{Provocation (激将法)} & “你是不是玩不起？/你这就是怂了。” & \textbf{劫持 Se 行动力}：利用 SP 渴望被认可为“强者/有趣”的心态，\tabularnewline
ESTP & \textbf{Ni 劣势} & \& \textbf{捧杀} & “只有你这种高手才能帮我。” & 通过质疑其能力或胆量，诱导他们做出冲动的承诺或牺牲。\tabularnewline
\midrule
\textbf{ISFP} & \textbf{Te 劣势} & \textbf{Intellectual Bullying (智力霸凌)} & “你的逻辑完全不通，听我的就行。” & \textbf{碾压 Te 话语权}：利用 SP 不善言辞辩论的弱点，\tabularnewline
ESFP & \textbf{Ni 劣势} & \& \textbf{宿命论暗示} & “我们是命中注定的，你逃不掉。” & 用复杂的术语或歪理让其闭嘴，确信自己“脑子不好”而交出控制权。\tabularnewline
\bottomrule
\end{tabular}}
\end{table}


\begin{table}[h!]
\caption[反PUA防御矩阵：如何利用认知功能进行框架重置]{反PUA防御矩阵：如何利用认知功能进行框架重置。
核心心法：PUA 的本质是试图把你拉进你的劣势功能(Grip)里打败你。防御的唯一方式是拒绝进入那个战场，强制把战场拉回你的主导/辅助功能领域。}
\centering
\resizebox{\textwidth}{!}{
\begin{tabular}{ccccc}
\toprule
防御者类型 & 敌方主攻点 & 防御武器 (Function) & 战术逻辑 (The Logic) & 核心反杀话术 (Counter-Script)\tabularnewline
\midrule
\multicolumn{5}{c}{NT 分析师：用“客观标准”粉碎“主观打压”}\tabularnewline
\midrule
\rowcolor{intjpurple!20} \textbf{INTJ} & \textbf{Se 焦虑} & \textbf{Te Parent} & \textbf{现实审计 (Reality Audit)} & “你的评价很有趣，但在这个领域，\tabularnewline
\rowcolor{intjpurple!20} ENTJ & (你真无趣/死板) & \textbf{外倾思考} & 拒绝自证。用客观标准替代对方的 & 我的KPI/成果是业内顶尖的。”\tabularnewline
\rowcolor{intjpurple!20} & & & 评价体系。把“人身攻击”变成“数据讨论”。 & “我不需要显得有趣，我只需要有效。”\tabularnewline
\midrule
\textbf{INTP} & \textbf{Fe 恐慌} & \textbf{Ti Hero} & \textbf{定义拆解 (Deconstruction)} & “你刚才撤回关心的行为，逻辑上\tabularnewline
ENTP & (撤回关注/冷暴力) & \textbf{内倾思考} & 无视情绪波动。像解剖尸体一样 & 属于操纵性惩罚。这对解决问题无效。”\tabularnewline
& & & 分析对方行为的逻辑矛盾。 & “如果你无法清晰表达需求，沟通终止。”\tabularnewline
\midrule
\multicolumn{5}{c}{NF 外交官：用“事实记录”对抗“记忆篡改”}\tabularnewline
\midrule
\textbf{INFJ} & \textbf{Se 混乱} & \textbf{Ti Child} & \textbf{证据留痕 (Documentation)} & “我查了聊天记录(数据)，事实并非如你所说。”\tabularnewline
ENFJ & (煤气灯/你记错了) & \textbf{内倾思考} & 相信逻辑记录，不相信模糊的记忆。 & “你的感受是你的，但物理事实是唯一的。”\tabularnewline
& & & 建立外部的“真相锚点”(如日记/截图)。 & “我不接受你对我记忆的修改。”\tabularnewline
\midrule
\textbf{INFP} & \textbf{Te 无能} & \textbf{Si Child} & \textbf{经验回溯 (Reference)} & “我过去独自解决过比这更难的事。”\tabularnewline
ENFP & (没有我你不行) & \textbf{内倾感觉} & 调取过去的成功经验。 & “我有手有脚，也有存款。分开是痛苦的，”\tabularnewline
& & & 证明“独立生存”的历史事实。 & “但绝不是致命的。”\tabularnewline
\midrule
\multicolumn{5}{c}{SJ 守护者：用“止损核算”对抗“恐惧画饼”}\tabularnewline
\midrule
\textbf{ISTJ} & \textbf{Ne 恐惧} & \textbf{Te Parent} & \textbf{风险评估 (Risk Assessment)} & “我计算了一下，留下来被你消耗的成本，”\tabularnewline
ESTJ & (离开会很惨) & \textbf{外倾思考} & 将“未知的恐惧”量化为“具体的损益表”。 & “远高于离开后重组生活的成本。”\tabularnewline
& & & 用利益计算压倒情绪恐吓。 & “不画饼，看报表。”\tabularnewline
\midrule
\textbf{ISFJ} & \textbf{Ti 自卑} & \textbf{Ti Child} & \textbf{逻辑自洽 (Consistency Check)} & “你说爱我，但行为上却在伤害我。”\tabularnewline
ESFJ & (为了赢回爱) & \textbf{内倾思考} & 强制检查言行一致性。 & “这在逻辑上是自相矛盾的。”\tabularnewline
& & & 只要发现一个矛盾点，就全盘否定。 & “我不玩这个竞争游戏，我退出。”\tabularnewline
\midrule
\multicolumn{5}{c}{SP 探索者：用“价值/逻辑”对抗“激将捧杀”}\tabularnewline
\midrule
\textbf{ISTP} & \textbf{Fe 羞耻} & \textbf{Ni Child} & \textbf{动机洞察 (Intent Analysis)} & “你试图通过激将法让我免费干活？”\tabularnewline
ESTP & (你敢不敢/是不是怂) & \textbf{内倾直觉} & 直接点破对方背后的意图。 & “这个策略对我太透明了，换一招吧。”\tabularnewline
& & & 预判后果，拒绝冲动。 & “这不仅仅是怂，这是愚蠢的冒险。”\tabularnewline
\midrule
\textbf{ISFP} & \textbf{Te 霸凌} & \textbf{Fi Hero} & \textbf{价值隔离 (Value Shield)} & “你说的道理也许都对，但我不喜欢。”\tabularnewline
ESFP & (你脑子笨听我的) & \textbf{内倾情感} & 承认逻辑弱势，但强调价值否决权。 & “我不需要辩赢你，我只需要拒绝你。”\tabularnewline
& & & “千金难买我乐意/不乐意”。 & “你的逻辑无法解释我的感受，所以无效。”\tabularnewline
\bottomrule
\end{tabular}}
\end{table}


\begin{table}[h!]
\caption[关系齿轮表：基于认知功能的深度兼容性分析]{关系齿轮表：基于认知功能的深度兼容性分析。
核心原理：最好的关系不是“相似”(会竞争)，也不是“相反”(会无法沟通)，而是“互补”(齿轮咬合)。}
\centering
\resizebox{\textwidth}{!}{
\begin{tabular}{cccc}
\toprule
关系类型 & 组合示例 & 动力学机制 (Mechanics) & 体验描述\tabularnewline
\midrule
\textbf{黄金配对} & \textbf{INTJ + ENFP} & \textbf{英雄救美 (Hero $\leftrightarrow$ Inferior)} & \textbf{极度吸引与治愈}。\tabularnewline
(Golden Pair) & & INTJ 的 Ni Hero 给 ENFP 混乱的未来提供方向(Si Inferior)。 & 就像找到了失散多年的另一半。\tabularnewline
& & ENFP 的 Ne Hero 给 INTJ 紧绷的当下带来无限可能(Se Inferior)。 & 甚至能无意识地激活彼此的潜意识。\tabularnewline
\midrule
\textbf{照镜子} & \textbf{INTJ + INTJ} & \textbf{共鸣与死局 (Hero $\leftrightarrow$ Hero)} & \textbf{懂你，但无法救你}。\tabularnewline
(Mirror) & & 互相秒懂对方的 Ni 愿景和 Te 效率。 & 容易陷入 Ni-Fi Loop 共振。一旦发生矛盾，\tabularnewline
& & 但两人都缺乏 Fe/Si，没人能处理细节和情绪。 & 谁也不愿先低头(Fi)，冷战到死。\tabularnewline
\midrule
\textbf{教师/学生} & \textbf{INTJ + INTP} & \textbf{家长引导 (Parent $\leftrightarrow$ Child)} & \textbf{智力巅峰，生活低能}。\tabularnewline
(Pedagogue) & & INTJ 的 Te Parent 极其欣赏 INTP 的 Ti Hero (逻辑)。 & 聊天能聊通宵，哲学探讨极其过瘾。\tabularnewline
& & INTP 的 Ne Parent 极其滋养 INTJ 的 Ni Child。 & 但两人都是 Se 盲点/劣势，一起出门可能会饿死/迷路。\tabularnewline
\midrule
\textbf{超我冲突} & \textbf{INTJ + ISFJ} & \textbf{互相听不懂 (Misunderstanding)} & \textbf{外星人对话}。\tabularnewline
(Superego) & & 你的主导(Ni)是她的恶魔(Ni Demon)。 & 你觉得她平庸琐碎，她觉得你傲慢危险。\tabularnewline
& & 她的主导(Si)是你的恶魔(Si Demon)。 & 除非极高阶的整合，否则通常是互相折磨。\tabularnewline
\bottomrule
\end{tabular}}
\end{table}

\clearpage
参考书目简介
\begin{itemize}

	\item Book~\cite{beebe} demonstrates the bond between the eight
		types of consciousness Jung named and the archetypal complexes
		that impart energy and purpose to our emotions, fatasies, and
		dreams.
		It contains four parts: (1) theretical contributions, including the eight
		function-attitudes, (2) type and the MBTI, (3) history of type, and (4)
		applications of type.

\end{itemize}


\begin{thebibliography}{1}
    \bibitem{beebe} John Beebe.
    \textit{Energies and Patterns in Psychological Type: The reservoir of consciousness}.
    Routledge, 2017.
\end{thebibliography}


\end{document}
